% P48 | 5.9 Mis-specified volatility | Result: pricing at 20 percent when the realised volatility differs
% guide: dev_notes/paragraph_guides/P48_result_pricing_at_20_percent_when_the_realised_v.txt
Because the boundary $b(t;T)$ decreases in $\sigma$, the $20\%$ boundary lies below the optimal one when the realised volatility is lower than $20\%$, so the buyer exercises too late.
When the realised volatility is higher the $20\%$ boundary lies above the optimal one, and the buyer exercises too early.
At a realised volatility of $10\%$ the mean stopping time is \MisLowTau\ against \MisLowTauOpt\ for the optimal rule, and the buyer loses \MisLowLoss\ relative to optimal exercise.
Waiting too long costs a modest amount because the option is still close to the money.
At $30\%$ the mean stopping time is \MisHighTau\ against \MisHighTauOpt, and the loss relative to optimal exercise is \MisHighLoss.
Exercising too early therefore costs more than exercising too late by a factor of more than two in this example.
The mean profit itself, \MisLowPnL\ at $10\%$ and \MisHighPnL\ at $30\%$, is driven mostly by the price mismatch, since the option is worth less or more than the buyer paid.
The differences relative to the optimal rule are estimated with a standard error of up to \RegretSE\ at $10{,}000$ paths, so the loss at $10\%$ is only about two standard errors, which the next section tightens with ten times more paths.
